\subsection*{The complete three-primary quartic class-two image}
\begin{statement}{E.117}{Proved}{Certified third-division fields, integral Johnson selectors, and the complete mixed image}
Put $R=\mathbb Z_3$, $S=\{2,3,\infty\}$, and retain the
plus quartic, every tangent and four-coset map of E.103--E.105,
and the actual class-two origin $S_3$.
Write $G=G_{\mathbb Q}$, $G_S=G_{\mathbb Q,S}$,
$V=T_3E_+$, $C=T_3J_C$, and $A=W\oplus j(C)$ with $j=j_*/4$.
Identify $C$ with $V^3$ by the exact E.103 maps $\Phi$ and $\Psi/4$,
and put $B=\Lambda^2 A$.  The actual full third-division fields
of $E_+$ and its quadratic twist $E_4$ have degree sixteen,
class number one, and the literal maps (E117.11)--(E117.15).
For each of $V,V\chi_2,V(1),V(1)\chi_2,U_{\rm CM}$,
(E117.17) evaluates every finite-coefficient cohomology group
and its continuous inverse limit over $G_S$.  Fix the actual cocycle $h$ representing
the E.103 generator of $H^1(G,V)=R$.

The full arithmetic extension with geometric free pro-three
kernel $N$ is the pullback of an extension over $G_S$.
The isomorphism and inverse are (E117.10) below.
The actual $S_3$, its class-two action, and $h$ descend through
this same map.  In particular the group of first coordinates
over $G_S$ is $R^{13}$, and every one of these coordinates
lifts to the actual intrinsic class-two extension over $G_S$.
With the specified coordinate choices the full geometric-
conjugacy space there is
\[
 \operatorname{Sect}_{G_S}(N/\gamma_3N)/N
       \simeq R^{13}\times\bigl(R^{78}\oplus(\mathbb Z/3)^{10}\bigr).
 \tag{E117.1}
\]
This is a homeomorphism with the actual inverse constructions
in the proof.  In particular every pure compact global input
$(0,r_1[h],r_2[h],r_3[h])$ lifts.

Here is the complete unrestricted image.  Express its input,
using the unchanged E.103 coefficient maps, as
\[
 [a]=[b]+[e_{\mathbf r}],\qquad
 b\in Z^1_{\rm cts}(G,W),\quad
 e_{\mathbf r}=j_{\mathbf r}h=\sum_{i=1}^3r_i j_i h,
 \quad\mathbf r=(r_1,r_2,r_3)\in R^3.
 \tag{E117.2}
\]
Any initial representative is aligned to this one by its actual
abelian coboundary and the BCH conjugation (E113.16).
At every rational prime $p\notin S$ choose an actual prolongation,
a tame pro-three generator $t_p$, and arithmetic Frobenius $F_p$,
with $F_pt_pF_p^{-1}=t_p^p$.
Use the corresponding Tate basis and put
\[
 \begin{gathered}
 v_p=b(t_p)\in W,\qquad
 u_p=(F_{E,p}-1)^{-1}h(F_p)\in V\otimes\mathbb Q_3,\\
 \overline u_p=u_p+V\in(V\otimes\mathbb Q_3/V)^{F_{E,p}},
 \qquad \operatorname{ord}(\overline u_p)=3^{m_p},\\
 c_p=v_p\wedge j_{\mathbf r}u_p\in
              (W\wedge j(C))\otimes\mathbb Q_3 .
 \end{gathered}
 \tag{E117.3}
\]
Here $m_p=0$ means that the class is zero.
The matrix $F_{E,p}-1$ is invertible over $\mathbb Q_3$.
Its class $\overline u_p$ is the actual local elliptic
Kummer class of $[h]$, with the explicit boundary map in
(E117.38).  Thus its phase is retained when an inverse lift
is constructed.  The integer $m_p$ is independent of the unit
choice of generator of $H^1(G,V)$.

The input (E117.2) has an intrinsic class-two lift exactly when
both of the following tests hold:
\[
 \begin{gathered}
 \text{all fifty-five peripheral tests (E108.4)--(E108.5)
       hold for $b$ at $\ell=3$,}\\
 c_p\in W\wedge j(C)\quad(p\notin S),\qquad
 c_p\longrightarrow0\quad\text{in its three-adic lattice.}
 \end{gathered}
 \tag{E117.4}
\]
There is a scalar form of the second test.  For a vector in
a finite free $R$-module let $\nu_3$ denote the largest power
of three dividing it, with $\nu_3(0)=\infty$.
If $\mathbf r\ne0$, put $t=\min_i\nu_3(r_i)$.
Replace $u_p$ by the representative of $\overline u_p$ whose
two coordinates lie in $\mathbb Z[1/3]\cap[0,1)$.
This changes $c_p$ by an integral restricted family.
Consequently (E117.4)'s mixed test is precisely
\[
 \nu_3(v_p)+t\ge m_p\quad(m_p>0),\qquad
 \nu_3(v_p)+t-m_p\longrightarrow+\infty
                         \quad\text{on $\{p:m_p>0\}$.}
 \tag{E117.5}
\]
For $\mathbf r=0$ the mixed test is empty.
All limits here mean that at each fixed coefficient exponent
only finitely many primes fail the indicated divisibility.
No closure is added to the denominator of the support quotient.

Every nonempty fibre over (E117.2) is the entire torsor
\[
 R^{66}\oplus(\mathbb Z/3)^{10}\oplus\mathcal R_{\mathcal P_{\mathbb Q}}(R)^6
 \oplus\mathcal R_{\{p:p\equiv1\ (4)\}}(R)^3
 \oplus\mathcal R_{\{p:p\equiv3\ (4),\ p\ne3\}}(R)^3,
 \tag{E117.6}
\]
with the actual coefficient projections, cocycle sections and
inverse ratio maps proved below.  No lower stabilizer is quotiented
out.  Formulas (E117.29), (E117.33), and (E117.37) give the
particular correction and its inverse, including the true
compact action and its E.109 inner corrections.

For $R=\mathbb Q_3$ the same theorem holds after scalar
extension; the ten displayed torsion factors map to zero through
their actual rational coboundaries.  In (E117.4) the mixed requirement is simply
$c_p\to0$ in the rational coefficient space, and the first
line uses E.108's rational test.
Formula (E117.5) retains only its limit requirement, allowing
integer-valued, possibly negative valuations of $v_p$ and
$\mathbf r$.
Every restricted module in (E117.6) is replaced by its rational
$c_0$ space, with one common bound on denominators.

The integral mixed test has nonzero obstructions.  The integers
$m_p$ are unbounded.  More precisely, for each $n\ge1$ there
are infinitely many primes satisfying
\[
 \begin{gathered}
 p\equiv1\pmod8,\qquad \nu_3(p-1)=n,\qquad
 2\in(\mathbb F_p^{\times})^{3^n},\qquad
 F_{E,p}\equiv1\pmod{3^n},\qquad m_p\ge n.
 \end{gathered}
 \tag{E117.7}
\]
They are the actual Chebotarev sets specified in the proof.
For each such prime the E.108 toric input $(p,1,1)$ lifts.
For every nonzero compact vector $\mathbf r$ and every one
of these primes with $n>\min_i\nu_3(r_i)$, the combined input
$((p,1,1),\mathbf r)$ does not lift.
Both of its separate toric and compact inputs do lift.
The resulting exclusion also holds in every geometric quotient
mapping to this actual class-two group.
This evaluates the three-primary class-two stage of E.23;
the maps to later intrinsic quotients retain their further
lifting conditions.
\end{statement}
\begin{proof}
\emph{The genuine unramified based action.}
Let $F=\mathbb Q(\zeta_8)$ and use exactly the cusp
$O=[\zeta_8:0:1]$, parameter $s=y/z$ and unit tangent of
E.105.  Over $\mathcal O_F[1/2]$ the plane quartic is smooth:
its partial derivatives are $4x^3,4y^3,4z^3$.
The twelve cusp sections are disjoint there.  Indeed their
parameters are the roots of $T^4+1$, whose discriminant is
a power of two, and sections in distinct coordinate lines
cannot meet on the quartic.
At $O$ the exact completed branch is
\[
       x/z=\zeta_8(1+s^4)^{1/4},\qquad y/z=s.
 \tag{E117.8}
\]
At every finite place not dividing $2\cdot3$ its coefficients
are integral, and the specified tangent is a unit.

Apply SGA1 XIII Corollary2.9 and construction2.10 to this
smooth proper curve and its relative normal-crossing divisor.
The construction identifies the prime-to-residue-characteristic
geometric cover categories of the generic and special fibres;
the last part of the proof of Corollary2.12 proves that this
specialization is an isomorphism.  Lemma2.11 identifies the
actual corresponding inertia subgroups
\cite[XIII,2.9--2.12 and their proofs]{E_SGA1}.
Here is the basepoint refinement used in this application.
Over the strict henselization at a good place, a tame branch
of a three-power cover at (E117.8) has the Kummer form
$z^{3^a}=s u(s)$ with $u(s)$ a unit, after an etale change.
Hensel's lemma supplies the roots of its leading unit and all
three-power roots of unity over that strict henselization.
Evaluating its Puiseux fibre at the fixed unit tangent thus
commutes with the specialization functor.
It gives an isomorphism of the cover categories with their
specified fibre functors, not just their unpointed categories.
Every based finite three-power quotient is consequently fixed
by arithmetic inertia.  Taking the inverse limit proves that
the actual tangential action on the full $N$ is trivial there.
The same conclusion for the closed normal subgroup generated
by these inertias follows from its being an action homomorphism.

Put $I=\ker(G\to G_S)$.  Since $F/\mathbb Q$ ramifies only
at two, $I\subset G_F$, and it is also the kernel of
$G_F\to G_{F,S_F}$.  To check the last assertion, the Galois
closure over $\mathbb Q$ of an extension of $F$ unramified
outside $S_F$ is still unramified outside $S$.
Thus the preceding action is trivial on $I$.

Let $\mathcal E$ be the actual arithmetic extension with
geometric kernel $N$, and let $s_F$ be the specified tangent
section on $G_F$.  The group $N$ has rank
$2\cdot3+12-1=17$ and centre one.
For clarity the latter assertion follows integrally from its
Magnus embedding in $1+(X_1,\ldots,X_{17})
\mathbb F_3\langle\!\langle X_1,\ldots,X_{17}\rangle\!\rangle$.
The embedding is faithful: evaluation in the group algebra
of any finite three-group quotient is defined because its
augmentation ideal is nilpotent, and these quotients separate
the free pro-three group.  A power series commuting with $X_i$
contains only powers of $X_i$, by comparing its first letter
and moving initial $X_i$'s to the end.  Commuting with both
$X_1$ and $X_2$ therefore makes it constant.  Its augmentation
is one, proving the assertion.

Both $e s_F(i)e^{-1}$ and $s_F(gig^{-1})$, for $e$ above $g$,
centralize $N$.  They have the same arithmetic image, so their
ratio is in $N\cap C_{\mathcal E}(N)=Z(N)=1$.
It follows that the compact subgroup $s_F(I)$ is normal in
$\mathcal E$ and meets $N$ trivially.  Define
\[
 \mathcal E_S=\mathcal E/s_F(I),\qquad
 1\longrightarrow N\longrightarrow\mathcal E_S
      \longrightarrow G_S\longrightarrow1.
 \tag{E117.9}
\]
The exact pullback and inverse are
\[
 \begin{gathered}
 \mathcal E\longrightarrow\mathcal E_S\times_{G_S}G,
          \qquad e\longmapsto([e],\pi(e)),\\
 ([e],g)\longmapsto e\,s_F(\pi(e)^{-1}g).
 \end{gathered}
 \tag{E117.10}
\]
The second expression is independent of the lift $e$;
its displayed argument lies in $I$.
These are continuous inverse homomorphisms, as one checks
by the normality identity just proved.

Every construction of E.103 and E.105 commutes with this
quotient.  In particular the chosen four coset lifts map to
their actual images, their factors $h_i(g)$ map to the same
factors in $G_{F,S_F}$, and the formula
$B(e)=\sum_i e e_i s_F(h_i(g))^{-1}e_{g i}^{-1}$ maps
term by term to the identical sum in the quotient.
The roots of two occurring in the frame cochains are
unramified outside $2,3$.  Thus the subsequent toric
correction and the central four-coset construction also
descend.  This proves the assertion for the specified
$S_3$, its actual $D_g$ and its linear action.
Changes of coherent roots or coset lifts remain the actual
fixed geometric conjugations already specified in E.105.
They preserve this descent and are not set equal to zero.

\emph{The two actual third-division fields and all local maps.}
For $A=1,4$ let $E_A:y^2=x^3+Ax$, choose the positive real root
$a_A$ of the polynomial below, and put
\[
 \begin{gathered}
 P_A(T)=T^8+216A^3T^4-3888A^6,
 \qquad L_A=\mathbb Q(i,a_A),\\
 Y_A=a_A/3,\quad
 X_A=\frac{12A^2a_A^2}{a_A^4+108A^3},\quad
 P=(X_A,Y_A),\quad IP=(-X_A,iY_A),\\
 r_A=\frac{3X_A^2+3A}{2A},\qquad
 \zeta_3=\frac{-1+ir_A}{2}.
 \end{gathered}
 \tag{E117.11}
\]
The exact field arithmetic verifies $r_A^2=3$, both curve
equations, $[2]P=-P$, $[3]P=[3]IP=0$, and all nine distinct
points $uP+vIP$, $u,v\in\mathbb F_3$.  These identities can
also be checked directly from the elliptic addition formulas
and $P_A(a_A)=0$.  The duplication equation $[2]P=-P$ is
equivalent to $3X_A^4+6AX_A^2-A^2=0$ together with its
ordinate equation.  Since $T^2+1$ is irreducible over
$\mathbb F_3$, the two points form a basis.  Consequently
every point of $E_A[3]$ is defined over $L_A$.  Conversely
$a_A=3Y_A$ and $i=Y(IP)/Y(P)$ belong to the field of their
coordinates.  Thus, with both inclusions specified,
\[
 L_A=\mathbb Q(E_A[3]),\qquad [L_A:\mathbb Q]=16,
 \qquad \mu_3\subset L_A.
 \tag{E117.12}
\]

Here are the exact degree and class certificates used in the
argument.  The reduced defining polynomials for $A=1,4$ are
respectively
\[
 \begin{split}
 Q_1(T)={}&T^{16}+6T^{12}+39T^8-18T^4+9,\\
 Q_4(T)={}&T^{16}-4T^{15}+14T^{14}-36T^{13}+80T^{12}
 -152T^{11}+256T^{10}-372T^9\\
 &+451T^8-444T^7+328T^6-172T^5+80T^4-48T^3
 +26T^2-8T+1.
 \end{split}
 \tag{E117.13}
\]
The program constructs the compositum of $P_A(T)$ and
$T^2+1$, proves the returned polynomial irreducible, and
retains the actual radical inclusions in it.  If its primitive
element is $i+k a_A$, the stored reduction polynomial gives
this element in $\mathbb Q[T]/Q_A$.  The inverse power-basis
matrix gives the inverse primitive polynomial.  Both matrix
products and both substituted polynomial identities are checked.
All elliptic and automorphism calculations therefore occur in
these specified fields, with maps in both directions.
The full class-group certificate, using PARI's default
\texttt{bnfcertify} with flag zero
\cite[General number fields, \texttt{bnfcertify}]{E117_PARI}, is
\[
 \begin{array}{c|c|c|c|c}
 A&\operatorname{disc}(L_A)&h(L_A)&2:\ (e,f)&3:\ (e,f)\\\hline
 1&2^{36}3^{14}&1&(8,2)&(8,2)\\
 4&2^{32}3^{14}&1&(8,2)&(8,2).
 \end{array}
 \tag{E117.14}
\]
There is exactly one prime in each of the last two entries.
The discriminants also prove unramifiedness outside $2,3$.
The actual BNF objects, unconditional certification returns,
prime decompositions, inverse field maps and textual records
are retained with the certificate.

The field automorphisms are literal substitutions:
$\tau(i)=i$, $\tau(P)=P+IP$, and $c(i)=-i$, $c(P)=P$.
The formula for the primitive element just given converts
these prescriptions to actual polynomials modulo $Q_A$.
The program verifies all sixteen substitutions and
\[
 \tau^8=c^2=1,\quad c\tau c=\tau^3,\qquad
 \rho_1(\tau)=\begin{pmatrix}1&-1\\1&1\end{pmatrix},
 \quad\rho_1(c)=\begin{pmatrix}1&0\\0&-1\end{pmatrix}
 \quad\text{over }\mathbb F_3.
 \tag{E117.15}
\]
The residual Weil-basis map is explicit as well.  Put
$m=(3X_A^2+A)/(2Y_A)$.  The flex tangent functions at $P$
and $IP$ have the same leading pole coefficient at infinity,
so their divisor-pairing formula gives
\[
 e_3(P,IP)=-\frac{(i-1)Y_A+2mX_A}{(1-i)Y_A+2imX_A}
              =\zeta_3^2.
 \tag{E117.16}
\]
The last equality is verified in both actual number fields.
This retains the sign of the residual Tate generator.
The eight powers of $\tau$ and their products with $c$ are
distinct.  They give the entire Galois group
$C_8\rtimes_{3}C_2$.  Both generators take the actual
$\zeta_3$ of (E117.11) to its inverse, exactly their determinant
action.  The single primes in (E117.14) imply that the
decomposition group at each of $2,3$ is this entire group.
The twist $V\chi_2\simeq T_3E_4$ is the actual map
$(x,y)\mapsto(2x,2\sqrt2\,y)$ and its inverse, as in E.112.

\emph{Finite-base cohomology and its integral inverse coordinates.}
Retain $U=U_{\rm CM}$ with the row, section and kernel of
(E108.24).  For all five lattices
$N\in\{V,V\chi_2,V(1),V(1)\chi_2,U\}$, the calculation is
\[
 \begin{gathered}
 H^0(G_S,N/3^n)=H^2(G_S,N/3^n)=0,
 \qquad H^1(G_S,N/3^n)\simeq\mathbb Z/3^n,\\
 H^1_{\rm cts}(G_S,N)=R,\qquad
 H^2_{\rm cts}(G_S,N)=0.
 \end{gathered}
 \tag{E117.17}
\]
Here is the complete finite-base argument.  Restrict first to
the appropriate $L_A$.  Its degree sixteen is a unit in $R$;
averaging the cochain action and corestriction divided by
sixteen give the inverse descent maps.  Every residual
coefficient is trivial over $L_A$, as are $\mu_3$.
For a chosen primitive $\zeta_3$, the coefficient isomorphism
is $(N/3)\simeq\mu_3\otimes(N/3)(-1)$.  The Kummer
sequence and (E117.14) identify degree two over
$\mathcal O_{L_A}[1/6]$ with its Brauer group tensored with
$(N/3)(-1)$.  Its Brauer group of exponent three is the
augmentation kernel in the permutation module on the two
places above $2,3$; there are no real places.  Each place
is fixed by the whole Galois group.  The invariant spaces
of $(N/3)(-1)$ are zero.  Explicitly over $\mathbb F_9$ let
$\alpha$ be the character of $\tau$ of order eight.
The exponents for $V$, its determinant twist, and $U$ are
respectively
$(1,3)$, $(5,7)$, and $(2,6)$; a Tate twist by $-1$
adds four.  None contains exponent zero.  The matrix
certificate independently stacks the actual equations
$(\rho(g)-1)x=0$ for all sixteen elements, both before and
after this twist, and obtains rank two in every case.
The same calculation applies locally because both
decomposition groups are full.  This proves $H^0=H^2=0$
at exponent one with the precise Tate twist retained.
The Kummer and Brauer maps are the localization maps of
\cite[II.0.1(f), II.1.7(b), II.2.9]{E_MilneADT}.

Induct on $n$ in the actual short exact coefficient sequences.
The mod-three vanishing gives $H^0=H^2=0$ at every exponent
and surjectivity of degree-one transitions.  Complex
conjugation has one plus and one minus dimension on each
of the five lattices.  The finite global Euler characteristic
therefore gives $|H^1(G_S,N/3^n)|=3^n$
\cite[I.4.15; I.5.1 and proof]{E_MilneADT}.  A nonzero initial
class lifts successively; Nakayama, or the coefficient exact
sequence, makes these groups cyclic with compatible
generators.  Their finite cocycles can be aligned by
coboundaries, since $H^0=0$.  This constructs an actual
continuous cocycle $h_N$ and the inverse map reading its
unique compatible scalars modulo $3^n$.  The finite groups
are Mittag--Leffler, so their derived limits prove (E117.17)
with the continuous-cochain comparison
\cite[Proposition1.6, Theorem2.2]{E_JannsenContinuous}.
For $N=V$ choose the initial reduction of the fixed E.103
generator $h$; inverse inflation and the cyclic calculation
show that it is nonzero.  The subsequent representative
alignment is by actual coboundaries, so it retains precisely
that generator and its stated marking.

For the first four lattices, tame inertia at a good prime
outside $S$ lies in the kernel of $F_N-p$.  Its eigenvalues
have complex absolute values $\sqrt p$ or $p\sqrt p$;
this kernel is zero.  Thus their unrestricted $H^1$ is
the same $R$, with inverse inflation.  For $U$, the fixed
space of $p^{-1}F_U$ is zero at split primes and a rank-one
direct summand at inert primes.  To see both assertions
integrally, pass faithfully to the unramified quadratic
coefficient ring where the CM characters are $\alpha,\beta$.
At a split prime their ratio is not one: equal eigenvalues
would give an integer elliptic trace of square $4p$.
At an inert prime $F_V^2=-p$, so $(p^{-1}F_U)^2=1$ and
the two eigenspaces have rank one.  Their integral
projectors are $(1\pm p^{-1}F_U)/2$ and descend to $R$.
In the real basis $e^2-f^2,ef$ an antidiagonal CM-normalizer
matrix on $V$ has the form
$\left(\begin{smallmatrix}a&b\\b&-a\end{smallmatrix}\right)$,
with $p=-a^2-b^2$.  Its actual plus projector is
\[
 P_p^+=\frac12\left[
 I_2+\frac1p
 \begin{pmatrix}a^2-b^2&ab\\4ab&b^2-a^2\end{pmatrix}
 \right].
 \tag{E117.18}
\]
Choose the first diagonal entry of $P_p^+$ that is a unit;
one exists since its trace is one.  Its column $w_p$ and
the corresponding row divided by that diagonal entry
give $\ell_pw_p=1$, $w_p\ell_p=P_p^+$.
These are the actual residue section and inverse.
Finite localization and (E117.17) lift every finite-support
residue, and finite unramified cohomology gives compatible
lifts by compactness.  Hence
\[
 H^1(G,U)=R\oplus
 \mathcal R_{\{p:p\equiv3\ (4),\ p\ne3\}}(R).
 \tag{E117.19}
\]
The map reads the $\ell_p$-coordinates of actual tame
inertia and then the scalar of the remaining unramified
cocycle; its inverse uses the chosen compatible residue
lifts.  More explicitly choose one integral lifted cocycle
for each residue basis column $w_p$.  Each is uniformly
bounded in its integral coefficient lattice.  For a restricted
family $(t_p)$ the sum of $t_p$ times these cocycles converges
uniformly, since $t_p\to0$; it is a continuous cocycle with
the prescribed residues.  Subtraction leaves the single
unramified scalar.  These two operations are continuous
inverse maps.  The omission of $3$ records its inclusion in $S$.

Finally apply the actual character matrices and inverse
CM maps of (E108.20)--(E108.24).  They decompose the central
coefficient into $\Lambda^2W$, six Tate lines, three
$R(1)\chi_{-1}$ lines, fifteen $V(1)$, eighteen
$V(1)\chi_2$, and three $U$ factors.
E.108 gives $H^1(\Lambda^2W)=R^{30}\oplus(\mathbb Z/3)^{10}$;
each torsion column is literally
$g\mapsto(\chi_3(g)^2-1)/3$ in its trivial character line.
The six Tate lines have all rational-prime valuation
coordinates, and the three twisted Tate lines have the
Gaussian ratios $\varpi_p/\bar\varpi_p$ for $p\equiv1\pmod4$.
Together with (E117.17)--(E117.19) this proves (E117.6),
including all coefficient maps and their inverses.

\emph{The literal compact primitives over the finite base.}
At $p\notin S$, a $V$-cocycle has inertia vector in
$\ker(F_{E,p}-p)$.  Its determinant is
$p\#E_+(\mathbb F_p)\ne0$, so this vector is zero.
The cocycle $h$ itself therefore factors through $G_S$.
This is the actual inverse inflation map just proved in (E117.17).

Put $B_W=\Lambda^2W$, $M=W\wedge j(C)$, and retain
$B_{\rm T}$ and $U_{\rm CM}^3$ from (E113.20).
E.108 and (E117.17) prove at every finite coefficient exponent
\[
 H^2(G_S,(B_W\oplus M\oplus U_{\rm CM}^3)/3^a)=0,
 \quad H^3(G_S,-)=0,
 \quad H^1(G_S,-)\text{ finite}.
 \tag{E117.20}
\]
For $B_W$ this is the proved $K_2(\mathcal O_{F,S_F})/3^a=0$
comparison, followed by restriction/corestriction of degree four.
For $M$ the exact coefficient isomorphism is
$V(1)^{15}\oplus(V(1)\chi_2)^{18}$ from E.108.
The other three summands are the $U$ of (E117.17).
The finite degree-one groups kill the derived inverse-limit
terms, so the continuous second groups in (E117.20) vanish.

We also need an injective map, rather than a vanishing
assertion, for the nine Tate lines:
\[
 H^2_{\rm cts}(G_S,B_{\rm T})
          \lhook\joinrel\longrightarrow H^2_{\rm cts}(G,B_{\rm T}).
 \tag{E117.21}
\]
Restrict to $F$; the nine characters become trivial and
restriction is injective by corestriction and division by four.
The class number of $F$ is one by E.103, so the Kummer
sequence on $\mathcal O_{F,S_F}$ has no Picard-group kernel.
Its degree-two map to $F$ is the injective map on Brauer
groups.  This proves injectivity at every finite exponent.
The finite $S_F$-unit groups modulo powers make $H^1$
finite and hence Mittag--Leffler.  On the unrestricted field
the Kummer $H^1$ transitions are surjective.
The same continuous-cochain comparison used in E.113 now
proves (E117.21), including its inverse-limit assertion.

The E.113 Weil cup therefore has a primitive $k_S$ on $G_S$.
For each $i$ the actual two-cocycle
$D_g j_i h(t)$ has zero class in every remaining summand
by (E117.20), and in the Tate lines by (E113.21) and
(E117.21).  Fix actual continuous cochains $p_i$ on $G_S$
with $dp_i(g,t)=D_gj_i h(t)$.
Here and below a primitive can be constructed by finite
linear equations: at each exponent its possible cochains
are a nonempty finite torsor under the finite group of
one-cocycles.  The inverse systems have the finite
intersection property, since any finite set of compatibility
requirements can be met at its largest exponent.
Their compact inverse limit supplies compatible cochains.
Fix orders on the finite sets and choose a compatible branch.
All arithmetic actions in these equations are the descended
actual actions just proved.

Define, with $e^{-1}:R(1)\to\Lambda^2V$ the actual inverse
Weil pairing,
\[
 \begin{gathered}
 \nu_{\mathbf r}=\tfrac12(\Lambda^2j_{\mathbf r})e^{-1}k_S,
 \qquad t_{\mathbf r}=\sum_i r_i p_i+\nu_{\mathbf r},\\
 dt_{\mathbf r}=\Omega_{e_{\mathbf r}},\qquad
 \Omega_a(g,t)=D_ga(t)+\tfrac12a(g)\wedge g a(t).
 \end{gathered}
 \tag{E117.22}
\]
This is a polynomial continuous choice in the three compact
scalars.  Changing from an earlier E.113 primitive $k$ to
$k_S$ changes the coordinates by the displayed central
one-cocycle involving $k-k_S$; no origin is changed.

\emph{Integral polarization and the entire actual compact action.}
Put $H=C=T_3J_C$.
The rational elliptic quotient maps of E.103 satisfy
$\Phi\Psi=\Psi\Phi=[4]$, and $\Psi$ is the polarization
dual of $\Phi$.  Therefore the actual Tate maps
\[
 \phi=\Phi/2:H\longrightarrow V^3,\qquad
 \psi=\Psi/2:V^3\longrightarrow H
 \tag{E117.23}
\]
are inverse symplectic isomorphisms over $R$.
For a concrete common CM basis choose an $R[i]$ generator
$t$ of $V$ lifting the real point $P$ in (E117.11).
Write $c(t)=u t$.  Then $u\equiv1\pmod3$ and
$u\bar u=1$.  Replacing $t$ by $u^{1/2}t$, with its unique
square root in $1+3R[i]$, makes $c(t)=t$ and preserves its
reduction.  Use $(t,It)$ and the actual Tate generator
$e(t,It)$, whose reduction is $\zeta_3^2$ by (E117.16).
Relative to any earlier coherent Tate generator, its unit
ratio is the actual Weil-pairing ratio, and the inverse
coordinate multiplies by the reciprocal unit.  Thus the
basis and its Tate twist have a specified transport at
every level; no independent root phase is assigned.
Relative to the original E.103 compact coordinate
$x_{\rm old}=\Phi x$, this is $x_{\rm new}=x_{\rm old}/2$,
with inverse multiplication by two.  The affine compact
inclusion is consequently $j_{\rm new}=2j_{\rm old}$.
Each old scalar $r_i$ in the statement is thus represented
by $r_i/2$ in the new basis.  On class-two coordinates the
literal morphism is $(x,z)\mapsto(Lx,\Lambda^2Lz)$,
where $L$ is this coefficient change and is the identity
on $W$.  Its inverse uses $L^{-1}$.  The actual off-diagonal
action changes to $\Lambda^2L\,D_gL^{-1}$; the affine
filling term changes by the same domain and target maps.
This is the stated coordinate transport, without any
identification of different arithmetic actions.

In this symplectic basis write
$H=\bigoplus_{i=1}^3(Re_i\oplus Rf_i)$ and
$q=\sum_i e_i\wedge f_i$.  The first factor is a rational
elliptic projection, so the coefficient row
$\varepsilon_1:\Lambda^2H\to R(1)$ extracting
$e_1\wedge f_1$ is equivariant and
$\varepsilon_1(q)=1$.  Thus the exact integral section of
$Z=\Lambda^2H/Rq$ is
\[
 s_1([z])=z-q\varepsilon_1(z),\qquad
 Z\xrightarrow{s_1}\Lambda^2H
       \xrightarrow{\text{quotient}}Z=\operatorname{id}.
 \tag{E117.24}
\]
It is independent of the representative $z$ and is equivariant.
Over $\mathbb Q_3$ its difference from the total-contraction
section in E.115 is exactly
$q(\langle z\rangle/3-\varepsilon_1(z))$.
This displays the comparison even though that total formula
need not preserve the three-adic lattice.

Here is an integral selector for the actual Johnson class.
Use the symplectic map
$f:\Lambda^3H(-1)\to\operatorname{Hom}(H,Z)$ of (E109.3)
and the actual inner map $\operatorname{ad}:H\to\operatorname{Hom}(H,Z)$.
The same integer matrices define these maps over $R$.
Their retained minors have determinants $\pm1$:
the twenty Johnson rows, in zero-based column-first Hom
coordinates, are
\[
 \begin{aligned}
 \mathcal I_f={}&(0,1,2,3,4,9,10,11,12,13,14,\\
 &23,24,25,26,27,28,38,39,42),\\
 \mathcal I_{\rm ad}={}&(0,1,2,3,4,14).
 \end{aligned}
 \tag{E117.25}
\]
The inverses of these actual minors, placed in the indicated
columns, give integer left inverses.  In particular both
images are saturated.  Contraction on $q\wedge H$ is
twice the identity, and $f(q\wedge h)=-\operatorname{ad}_h$.
These identities also prove that the induced map
$\bar f:L/(q\wedge H)\to\mathcal D$ is injective with
saturated image, where $L=\Lambda^3H(-1)$ and
$\mathcal D=\operatorname{Hom}(H,Z)/\operatorname{ad}H$:
the full Johnson image contains the inner image, so its
quotient has the same torsion-free complement.

Over $R'=R[i]$, set
$e_i^+=e_i-if_i$, $f_i^-=(e_i+if_i)/(2i)$.
Their basis block is
$\left(\begin{smallmatrix}1&-i/2\\-i&1/2\end{smallmatrix}\right)$,
of determinant one.  For a derivation $D$ into $Z$, define
$\Pi_3(D)$ by its two pure-CM coordinates
\[
 \bigl([e_2^+\wedge e_3^+]D(f_1^-),
       -[f_2^-\wedge f_3^-]D(e_1^+)\bigr).
 \tag{E117.26}
\]
The indicated coefficients do not meet $Rq$; this is
independent of every quotient lift.  Directly from the
three summands in $f$ one obtains
$\Pi_3 f=\pi_M$, $\Pi_3\operatorname{ad}=0$.
The CM normalizer consists of common diagonal blocks and
common antidiagonal blocks; substitution verifies
equivariance with their actual determinant twist.
It also verifies descent under $i\mapsto-i$.
For an entirely rational integral formula let
$M\subset L$ have the basis
\[
 \begin{split}
 m_0={}&e_1e_2e_3-e_1f_2f_3-f_1e_2f_3-f_1f_2e_3,\\
 m_1={}&e_1e_2f_3+e_1f_2e_3+f_1e_2e_3-f_1f_2f_3,
 \end{split}
 \tag{E117.27}
\]
where products denote ordered exterior products and the
inverse Tate twist is retained.  List the fourteen target
pairs lexicographically, omitting $e_3\wedge f_3$, and list $D(e_1),D(f_1),\ldots,D(f_3)$
in that order.  In these $84$ zero-based coordinates, the
first row of $\Pi_3$ has entries $1/4$ at $11,12,24$ and
$-1/4$ at $27$; the second has $-1/4$ at $10$ and $1/4$
at $13,25,26$.  All other entries are zero.
The quotient matrix, the basis-change matrix and their
inverses are retained, so this is also a literal formula
on the original quotient coordinates.  It gives
$\Pi_3\bar f|_M=1$, with no denominator divisible by three.

We now apply these maps to the genuine compact action,
not just its coefficient representation.  The proof of
the integral Harris--Pulte comparison in Hain--Matsumoto,
Theorem3 and \S8, and the exact group-ring/Magnus comparison
in BLLS Lemmas2.20--2.22 and Proposition2.23, hold for this
odd prime with their factor two
\cite{E109_HainMatsumoto,E109_BisognoLiLittSrinivasan}.
The map is $x\mapsto x-1$ on the group ring, followed by
its alternating degree-two quotient; its inverse at class
two divides the symmetric square by two.
Thus the same based action convention as in E.105 gives
$2[c^C]_{\rm out}=\bar f_*\mu$ before pure-CM projection.
The sixteen geometric diagonal automorphisms of the quartic
and their actual cycle translations are those in E.109.
Their average is integral over $R$; its image on
$L/(q\wedge H)$ is exactly $M$, since the other character
columns are nontrivial.  The origin translation contributes
$q\wedge H$ and hence vanishes in this quotient.
Restriction to $F$ followed by corestriction divided by four
therefore gives the actual class identity
\[
 \beta_3=2\Pi_{3*}[c^C]_{\rm out},\qquad
 2[c^C]_{\rm out}=\bar f_*\beta_3.
 \tag{E117.28}
\]
This is a statement about the actual three-adic class.
It makes no numerical assignment to that class from the
five-adic finite shadow.  All maps just used, including
the integral selector, are now specified over its own
coefficient and arithmetic source.

Use the actual filling morphism
$(x,z)\mapsto(px,p_Bz+Jx)$ and the descended origin $S_3$.
Its literal identities are
$D^C_gp=p_BD_g+J\rho_A(g)-\rho_Z(g)J$ and
$c^C=p_Bc^Aj+Jj-g(Jj)$, exactly as in (E109.13).
Inflation on $H^1$ is injective, so (E117.28) holds over
the same $G_S$.  The saturated Hom quotient allows its
actual zero-cochain to be lifted integrally.  Its remaining
inner cocycle lies in $H$, and (E117.17) supplies its three
unique scalars.  Absorb its coboundary into the Hom
zero-cochain to obtain
\[
 2c^C_g=f(\beta_3(g))+(gm-m)
      +\operatorname{ad}_{\sum_i\lambda_i h_i(g)},
 \qquad h_i=j_i^C h,\quad \lambda_i\in R.
 \tag{E117.29}
\]
Here $j_i^C$ and the scalars are in the original E.103
coordinates, using (E117.23) when the displayed selector
is applied.  The constants $m,\lambda_i$ retain the
actual action; none is supplied by a graded matrix.

Evaluating $f$ on $M$ sends the opposite CM eigenspace
into its pure exterior square and kills the same eigenspace.
Thus $f(\beta_3(g))(g h_i(t))$ lies in the actual $U^3$
summand of $Z$.  The certificate gives its rank-six image,
section and inverse.  Equation (E117.17) supplies actual
cochains $n_i$ there with this differential.
For the Weil primitive $k_S$ fixed above set
$Q_h=-h\otimes h/2+k_S(e\otimes f-f\otimes e)/2$.
Direct expansion gives $dQ_h=h\otimes gh$.
Let $K_i$ be the wedge-quotient image of
$\sum_j\lambda_j(j_j^C\otimes j_i^C)Q_h$.
Evaluating (E117.29) on $g h_i(t)$ then proves
\[
 t_i^C=\tfrac12(n_i+m h_i+K_i),\qquad
 dt_i^C(g,t)=D^C_g h_i(t).
 \tag{E117.30}
\]
Finally replace each earlier primitive $p_i$ by
$p_i+(\Lambda^2j)s_1(t_i^C-p_Bp_i+Jj_i h)$.
The parenthesized difference is an actual one-cocycle,
by the filling identity, so this is a permitted change.
It has the exact result
\[
                 p_Bp_i-Jj_i h=t_i^C.
 \tag{E117.31}
\]
This proves all integral compact, affine and origin
transfers needed below.

\emph{The exact obstruction reduction and its inverse.}
Let $P_W:B\to B$ be the actual exterior projection onto
$B_W$, followed by its given inclusion in $B$.
For a toric cocycle $b$, E.105 gives
$(1-P_W)\Omega_b|_{G_F}=0$ as a cochain.
Push out the pulled-back extension over $bS_3$ to
$\ker P_W$.  Its section on $G_F$ has central coordinate
zero.  Apply the same four-coset construction to that section.
If $B_b(e)$ denotes its literal four-term sum, the resulting
section has central coordinate $-B_b((0,g))/4$.
Consequently
\[
 U_b(g)=B_b((0,g))/4,\qquad
        dU_b=(1-P_W)\Omega_b.
 \tag{E117.32}
\]
It is a continuous cochain given by four actual group products;
it depends continuously, and polynomially in BCH coordinates,
on the fixed cocycle $b$.

Expanding the two first-cocycle identities gives
$d(b\wedge e)=-b\wedge ge+e\wedge gb$.
Equations (E117.22) and (E117.32) therefore give the inverse
cochain maps
\[
 \begin{gathered}
 z\longmapsto
 \widetilde z=z+t_{\mathbf r}+\tfrac12 b\wedge e_{\mathbf r}+U_b,
 \qquad
 d\widetilde z=-P_W\Omega_b-b\wedge g e_{\mathbf r},\\
 z=\widetilde z-t_{\mathbf r}
                  -\tfrac12 b\wedge e_{\mathbf r}-U_b.
 \end{gathered}
 \tag{E117.33}
\]
The two terms on the right lie in the distinct specified
coefficient summands $B_W$ and $M$.
Restriction to $G_F$ of the first term is exactly the
fifty-five cochains of E.108.  Restriction/corestriction of
degree four proves both directions of its stated test.
For an initially changed abelian representative use the
literal $\Gamma_m$ of (E113.16); the transformed central
coordinate changes by $\Gamma_m$ and by the difference
of the displayed added cochains.  This is the full gauge
map and its inverse, without identifying different origins.

\emph{Finite support before the derived limit.}
For $M=W\wedge j(C)$, localization at a finite set of primes
outside $S$, at coefficient exponent $3^a$, has support
term $\bigoplus_p R\Gamma(\mathbb F_p,(M/3^a)(-1))[-1]$.
Take the filtered direct limit in the finite set first.
Purity identifies its two-term complexes with
\[
 [M/3^a\xrightarrow{\ d_p\ }M/3^a],\qquad
             d_p=p^{-1}F_{M,p}-1.
 \tag{E117.34}
\]
These are the actual etale/Galois and inertia maps
\cite[II.0.1(f), II.1.7(b), II.2.9]{E_MilneADT}.
Both direct-sum coefficient transitions are onto.
Their derived inverse limit is thus the literal complex
$[\mathcal R(M)\xrightarrow d\mathcal R(M)]$, where
$\mathcal R(M)=\{(x_p):x_p\in M,\ x_p\to0\}$.
The finite-base second and third groups vanish by (E117.20),
and its finite first groups kill the adjacent derived limit.
The continuous-cochain comparison consequently gives
\[
 H^2_{\rm cts}(G,M)
       \simeq\mathcal R(M)/d\mathcal R(M).
 \tag{E117.35}
\]
Here $d\mathcal R(M)$ is its image, with no closure.
These are precisely the inverse-limit maps of
\cite[Proposition1.6 and Theorem2.2]{E_JannsenContinuous}.

The eigenvalues of $p^{-1}F_{M,p}$ have complex absolute
value $\sqrt p$, since $W$ has eigenvalues $p$ times roots
of unity and $C$ has the elliptic eigenvalues.
Hence every $d_p$ is invertible over $\mathbb Q_3$.
This proves that the candidate support correction is unique.

\emph{The literal mixed localizer and its compact-point inverse.}
In the central extension for the ordered cup
$\omega(g,t)=b(g)\wedge g e_{\mathbf r}(t)$, use the law
$(z,g)(w,t)=(z+gw+\omega(g,t),gt)$.
The compact cocycle is zero on inertia, so the central-zero
lifts $t_p^a\mapsto(0,t_p^a)$ are a homomorphism.
Multiplying the three factors of $F_pt_pF_p^{-1}$ gives
the central defect
\[
 \epsilon_p=-p\,v_p\wedge e_{\mathbf r}(F_p),\qquad
 \epsilon_p\longmapsto\epsilon_p+(F_{M,p}-p)z_p
 \quad\text{under }t_p\mapsto(z_p,t_p).
 \tag{E117.36}
\]
For example the two nonzero cup contributions are
$+b(F_p)\wedge e(F_p)$ and
$-(b(F_p)+p v_p)\wedge e(F_p)$.
The support map divides the defect by $p$, as it compares
the value at $t_p^p$ with the chosen value at $t_p$.
Thus (E117.35) sends this cup to
$[(-v_p\wedge e_{\mathbf r}(F_p))_p]$ and its unique
correcting family is
\[
 z_p=d_p^{-1}\bigl(v_p\wedge e_{\mathbf r}(F_p)\bigr).
 \tag{E117.37}
\]
The input family is restricted: each reduction of $b$ gives
a finite extension ramified at only finitely many primes,
so $v_p\to0$, and $h(F_p)$ is integral and uniformly bounded.

Write $F_{W,p}=pP_p$, where $P_p$ is the actual permutation
operator on the peripheral quotient.  The first-cocycle tame
relation gives $P_pv_p=v_p$.  Therefore
$d_p(v_p\wedge j_{\mathbf r}u)
=v_p\wedge j_{\mathbf r}(F_{E,p}-1)u$.
Since $\det(F_{E,p}-1)=\#E_+(\mathbb F_p)\ne0$,
(E117.37) is exactly $c_p$ in (E117.3).
The short exact sequence $0\to V\to V\otimes\mathbb Q_3
\to V\otimes\mathbb Q_3/V\to0$ over the unramified
local source has the explicit boundary
\[
 \overline u\longmapsto
       [F_p\longmapsto(F_{E,p}-1)u],\qquad
 H^1_{\rm unr}(\mathbb Q_p,V)=V/(F_{E,p}-1)V.
 \tag{E117.38}
\]
Both invariants and first cohomology of the rational
two-dimensional coefficient vanish here because $F_{E,p}-1$
is invertible.  This proves the identification and its inverse.
Under good reduction it is the actual finite group
$E_+(\mathbb F_p)[3^\infty]$.
Thus $m_p\le\nu_3(\#E_+(\mathbb F_p))$.
It can be read at that finite coefficient exponent by applying
the actual Frobenius and the corresponding finite cocycle
$h_a$ constructed below to the quotient in (E117.38).
This retains the arithmetic Frobenius marking at every level.

Changing $h$ by $dh_0$ changes $u_p$ by the same integral
$h_0$; replacing its generator by a unit does not change
$m_p$.  Replacing $u_p$ by its displayed fractional
representative changes it by an integral vector, uniformly
bounded for all $p$.  The resulting change in $c_p$ is
therefore integral and restricted, since $v_p\to0$.
Changing a tame basis changes both tensor factors by the
corresponding inverse units.  An inertia translate of Frobenius
does not change $h(F_p)$, since $h$ is unramified.
These are the actual representative maps proving independence
of all tests.
In the fixed tensor lattice, valuation of a nonzero pure tensor
is the sum of the valuations of its factors.
The fractional representative has valuation exactly $-m_p$
when $m_p>0$, and is zero when $m_p=0$.
This proves (E117.5), including its convergence requirement.

If the tests hold, (E117.35) gives a global primitive with
the prescribed restricted inertia corrections.
Here is its compatible construction.  At each finite exponent
the corrections have finite support.  Exactness of localization
gives a primitive; the surjective degree-one residue map, whose
next term is the zero finite-base $H^2$ of (E117.20), adjusts
its inertia values to the prescribed ones.
The remaining choices are a finite torsor under unramified
one-cocycles.  All finite compatibility conditions can be met
at their largest exponent, so their compact inverse limit
gives an actual continuous primitive.
Use the same proved construction of E.108 for the peripheral
summand and then (E117.33) for its literal class-two inverse.
Conversely any lift supplies these corrections by its tame
relation, so both tests are necessary.

\emph{All fibres and the finite-base coordinates.}
The abelian module $A$ has no global invariants: at a good
Frobenius its weights are one and two, so none of its
eigenvalues is one.  Thus the aligning abelian conjugator
is unique.  The difference of two lifted central coordinates
is exactly a $B$-valued one-cocycle, and a geometric central
conjugator adds its ordinary coboundary.
This proves the full fibre assertion and inverse ratio map;
substitution of (E117.17)--(E117.19) gives (E117.6).

For specificity all finite-base input coordinates are units
with actual Kummer maps.  Set $\epsilon=1+\sqrt2$,
$\varpi=1-\zeta_8$, and $\pi_\pm=1\pm\sqrt{-2}$.
The two primes of $F$ over three have these generators and
norm $9$; $N\varpi=2$, $N\epsilon=1$.
In each of the first two factors of (E108.1), the four
independent $S_F$-unit coordinates are
$\epsilon,\varpi,\pi_+,\pi_-$.
The norm-one third factor has the two coordinates
$\epsilon,\pi_+/\pi_-$, since its norm matrix in the
two rational unit coordinates $2,3$ is
$\left(\begin{smallmatrix}0&1&0&0\\0&0&2&2\end{smallmatrix}\right)$.
Together with the three compact coordinates this gives the
literal $R^{13}$ in (E117.1).

The first $66$ free central coordinates and the ten order-three
coordinates are those already constructed above.
Each of its six Tate lines has the two $S$-unit coordinates
$2,3$.  Its three $\chi_{-1}$ Tate lines have no $S$-coordinate:
two is ramified and three is inert in $\mathbb Q(i)$, so the
anti-invariant valuation module at these primes is zero.
There are no inert-prime residue coordinates in $S$.
Thus the full central group over $G_S$ is
$R^{78}\oplus(\mathbb Z/3)^{10}$.
For the particular section, choose the thirteen fixed
input cocycles just displayed.  Its obstruction is a polynomial
of degree at most two in their coordinates.  Every coefficient
has a primitive: use (E117.20) in the remaining summands,
(E117.21) and E.113 in the Tate summands, and polarization
over $R$ to separate the coefficients.  Fix these finitely
many primitives.  Their polynomial sum is a continuous
particular section for every input.
Subtracting it and reading the $78$ free and ten order-three central coordinates gives
the inverse of (E117.1).
At finite coefficient level the relevant section sets are
finite by arithmetic cohomological finiteness and the central
extension sequences.  Their inverse limits and their compact
geometric quotients give exactly this topology.

For rational coefficients every continuous cochain on the
compact source has a common denominator bound.
Scalar extension therefore commutes with the cochain complex.
The preceding maps, tests and inverses rationalize literally.
A sequence tending to zero in $\mathbb Q_3$ is bounded, which
is precisely the denominator convention in the statement.

\emph{Actual mixed nonlifting and unbounded local orders.}
Let $L_n=\mathbb Q(E_+[3^n])$, and let $\Gamma_n$ be its
actual Tate image.  The closed image of $G_{\mathbb Q(i)}$
acts by units in the unramified quadratic ring $R[i]$.
The actual residual element $\tau^4$ in (E117.15) is $-I$.
Any lift lies in $-1+3R[i]$; its successive $3^k$th powers converge to $-I$.
Closedness proves that $-I$ belongs to the actual Tate image,
and hence is central in every $\Gamma_n$.
For a one-cocycle on $\Gamma_n$ with coefficients $V/3$,
commuting an element with this scalar shows directly that
the cocycle is a coboundary, since $-2$ is a unit.
This is the elementary Sah calculation, with its actual scalar.

The reduction $h_1$ is nonzero by (E117.17), so its restriction
to $\ker(G_S\to\Gamma_n)$ is nonzero by inflation--restriction.
The residual representation is irreducible: the actual operator
$\rho_1(\tau)^2=\left(\begin{smallmatrix}0&1\\-1&0\end{smallmatrix}\right)$
has irreducible polynomial $T^2+1$ over $\mathbb F_3$.  The image of $h_n$ on that kernel is
therefore onto modulo three, and Nakayama's lemma gives the
whole $V/3^n$.  Consequently the actual finite quotient
\[
 g\longmapsto(h_n(g),\rho_n(g))
 \quad\text{has image }(V/3^n)\rtimes\Gamma_n.
 \tag{E117.39}
\]
This also specifies its exact field and its translation kernel.

The three-Sylow subgroup of $\Gamma_n$ lies in the principal-unit
group of $(R/3^n)[i]$ and has exponent at most $3^{n-1}$.
It cannot have the cyclotomic quotient of order $3^n$;
hence $\mu_{3^{n+1}}$ is not contained in $L_n$.
Adjoining $F$, of degree a power of two, preserves that fact.
Over $L_nF$, the Kummer extension generated by $2^{1/3^n}$
has cyclic three-kernel with conjugation character $\chi$.
It has no common degree-three quotient with the next
cyclotomic extension, whose conjugation character is trivial.
The characters are distinct modulo three.
Thus there is an automorphism of their compositum fixing
$L_nF$ and $2^{1/3^n}$ but acting nontrivially on
$\mu_{3^{n+1}}$.

The additional extension over $L_n$ just used has three-primary
composition factors with characters $1$ and $\chi$, and a
prime-to-three part.  Its common quotient with the translation
extension (E117.39) is trivial: every nonzero quotient of the
latter translation module has the irreducible $V/3$ as a
composition factor, which has dimension two.
All these statements concern the actual normal field
extensions and their conjugation actions over $\mathbb Q$.
We may therefore prescribe simultaneously the preceding
automorphism and a primitive translation in (E117.39).
Chebotarev applied to its conjugacy class gives infinitely
many primes with all properties in (E117.7)
\cite[VIII,Theorem7.4 and proof]{E_MilneCFT}.
Primitivity is preserved under conjugation.
Since $F_{E,p}-1$ is divisible by $3^n$ and $h(F_p)$ is
primitive, (E117.3) gives $m_p\ge n$.

Finally the input $(p,1,1)$ is
$b=\kappa_p\sum_{i=0}^3e_{A_i}$.
This is a primitive rational Tate line in $W$.
Each of its E.108 peripheral cups is a scalar multiple of
$\kappa_2\cup\kappa_p$; the two self-cups vanish because
two is invertible.  At primes over $p$ the residue field is
$\mathbb F_p$, and the whole three-primary tame symbol is
zero because $2$ is a $3^n$th power and $\nu_3(p-1)=n$.
At the prime above two its residue-unit group is trivial.
There are no other supports.  Thus the toric input lifts by
E.108.  Its mixed inertia vector at $p$ is primitive, so
(E117.5) fails whenever $n>\min_i\nu_3(r_i)$.
This proves all stated nonlifting assertions through the
actual projection to the intrinsic class-two group.

The certificate checks the symplectic contraction and its
pure-CM image, both tensor/cochain signs, the full tame word
and correction map, the actual $S_F$ norm coordinates,
the residual irreducibility and character separation, and
the finite tensor-valuation and inverse-coordinate identities.
Its finite computations supplement the preceding source,
localization and infinite-limit proofs; no numerical values
for an uncomputed higher Frobenius word are substituted.
\end{proof}
